% Source: https://artofproblemsolving.com/community/c3h2261902_radishes_money_and_a_confusing_math_problem
% Stable identity: aops-topic-2261902-q01
\textbf{(a)} One pound of whole radishes gives $t$ pounds of tops and $1-t$ pounds of bottoms, which bring in
\[3t+2(1-t)=2+t \text{ dollars}.\]
Since $t<1$, this is less than $\$3$, and certainly less than the $\$5$ he would get for the pound of whole radishes. He loses $5-(2+t)=3-t$ dollars, between $\$2$ and $\$3$, on every pound he cuts up. The blunder is that $\$3+\$2$ is the price of a pound of tops \emph{plus} a pound of bottoms, which is more than a pound of radishes.

\textbf{(b)} With $t=\dfrac{1}{2}$, one pound of tops and one pound of bottoms come from $2$ pounds of radishes, worth $2\cdot5=\$10$ when sold whole. He receives only $3+2=\$5$, so he loses $\boxed{\$5}$. (In agreement with (a): $2$ pounds at a loss of $3-\dfrac{1}{2}=\$2.50$ per pound.) This is the explanation ``he sold two pounds of radishes for the price of one'' in the riddle.
